% =====================================================================
%  CHAPTER 6: 社会文化因素与健康
% =====================================================================
\chapter{社会文化因素与健康}
\setcounter{chapter}{6}
\chapterintro{
掌握文化对人群健康的影响及作用模式（\textbf{教学难点}）；熟悉文化的概念、构成与特征；理解教育、风俗习惯、宗教等因素对健康的影响；了解科技进步、道德、非主流文化与健康的关系。
}

% ----- 6.1 文化概述 -----
\section{文化概述}

\begin{defbox}
\textbf{文化（Culture）}：广义指人类在社会历史实践过程中所创造的物质财富和精神财富的总和；狭义指社会的意识形态以及与之相适应的制度和组织机构。社会医学主要关注\keyword{狭义文化}——包括思想观念、价值体系、道德规范、语言文字、科学技术、宗教、艺术等精神文化要素。
\end{defbox}

\subsection{文化的构成}

\begin{tcolorbox}[colback=light, colframe=main, boxrule=0.5pt, arc=4pt, title=\textbf{文化构成的三个层次}, breakable]
\begin{enumerate}[leftmargin=*]
    \item \imp{物质文化}——人类创造的物质产品和技术成果，如工具、建筑、医疗设备等。直接影响人类生活条件和健康水平
    \item \imp{制度文化}——社会规范、法律、政治制度、社会组织形式等。通过约束和调节人类行为间接影响健康
    \item \imp{精神文化}——思想观念、价值体系、宗教信仰、伦理道德、科学知识、艺术审美等。是文化的核心，深刻影响人们的健康观念和健康行为
\end{enumerate}
\end{tcolorbox}

\subsection{文化的特征与类型}

\begin{keybox}
\textbf{文化的特征：}
\begin{enumerate}[leftmargin=*]
    \item \imp{历史继承性}——文化是历史传承的产物，通过社会化过程代代相传
    \item \imp{社会共享性}——文化是社会群体共享的知识和信念体系
    \item \imp{适应性}——文化随社会环境变化而演变和调整
    \item \imp{民族性和区域性}——不同民族和地区具有不同的文化特性
    \item \imp{习得性}——文化不是遗传的，而是后天学习和获得的
\end{enumerate}

\textbf{文化的类型：}
\begin{itemize}[leftmargin=*]
    \item \imp{主流文化（Mainstream Culture）}：社会中占主导地位的文化
    \item \imp{亚文化（Subculture）}：某一群体特有的区别于主流文化的文化模式
    \item \imp{反文化（Counterculture）}：与主流文化价值观念对立的文化
    \item \imp{跨文化（Cross-culture）}：不同文化背景的人群间的文化比较
\end{itemize}
\end{keybox}

% ----- 6.2 文化对健康的影响 -----
\section{文化对人群健康的影响}

\subsection{文化影响健康的作用模式（教学难点*）}

\begin{keybox}
\imp{（教学难点*）}\textbf{文化影响健康的四大作用模式：}

\begin{enumerate}[leftmargin=*]
    \item \imp{文化决定健康观念模式}：不同文化下的疾病认知、病因解释和健康信念显著不同，影响人们求医行为和健康行为选择。例如：不同文化对"健康"和"疾病"的定义存在差异；不同文化对疼痛、不适的表达方式不同

    \item \imp{文化塑造健康行为模式}：饮食文化（西方高脂饮食 vs. 地中海饮食 vs. 东方饮食）、饮酒文化、吸烟文化等直接影响健康。文化价值观和规范通过社会化过程内化为个人的日常行为习惯

    \item \imp{文化影响卫生服务利用模式}：传统文化和民间医学影响人们对现代医疗和传统医疗的选择偏好。不同文化群体对精神卫生服务的利用存在显著差异（病耻感、对病因的文化解释等）

    \item \imp{文化调节心理应激反应模式}：文化决定应激源的评价方式（什么是"压力"）、应对策略的选择（求助、回避、情绪调节等）和社会支持的可获得性。集体主义文化与个人主义文化在应对方式上存在显著差异
\end{enumerate}
\end{keybox}

\subsection{文化影响健康的特点}

\begin{enumerate}[leftmargin=*]
    \item \imp{潜移默化性}——文化对健康的影响往往是不自觉的、渐进的、内在化的
    \item \imp{广泛渗透性}——文化影响渗透到日常生活的方方面面，涉及饮食、起居、求医等各个环节
    \item \imp{持久稳定性}——文化价值观念和传统习俗根深蒂固，难以短时间改变
    \item \imp{双重性}——文化对健康既可能有积极促进作用，也可能有消极损害作用
    \item \imp{群体差异性}——不同文化背景人群的健康状况和健康行为存在显著差异
\end{enumerate}

% ----- 6.3 教育与健康 -----
\section{教育对健康的影响}

\begin{keybox}
\textbf{教育影响健康的机制：}
\begin{enumerate}[leftmargin=*]
    \item \imp{提高健康素养}——教育使人们获取和处理健康信息的能力增强，能够理解健康知识和遵从医嘱
    \item \imp{促进健康行为}——受教育程度高的人更可能采取积极健康的行为（合理饮食、适量运动、戒烟限酒等）
    \item \imp{增强应对能力}——教育提高人们应对生活压力和健康问题的能力
    \item \imp{改善职业和经济地位}——教育→更好的工作→更高的收入→更好的生活条件和医疗保障
    \item \imp{减少风险行为}——教育程度增高，吸烟、吸毒、酗酒等危险行为的发生率下降
    \item \imp{增强社会参与}——受教育程度高的人更多地参与社会活动，拥有更广泛的社会网络和支持
    \item 母亲受教育水平对儿童健康的\imp{独立保护效应}尤为显著
\end{enumerate}
\end{keybox}

% ----- 6.4 科技进步与健康 -----
\section{科技进步对行为与健康的影响}

\begin{enumerate}[leftmargin=*]
    \item \imp{积极影响}：医学技术进步（诊断、治疗、预防手段革新）、健康信息技术（电子健康档案、远程医疗、健康APP）、公共卫生监测和应急能力提升
    \item \imp{消极影响}：屏幕久坐行为增加（"电子久坐"）、互联网成瘾、睡眠紊乱（蓝光暴露）、信息误导（伪健康信息传播）、隐私泄露和基因歧视风险
\end{enumerate}

% ----- 6.5 风俗习惯与健康 -----
\section{风俗习惯对健康的影响}

\begin{enumerate}[leftmargin=*]
    \item \imp{有益健康的风俗}：中国传统饮食文化（均衡搭配、食疗养生）、饮茶习惯（抗氧化）、传统体育运动（太极拳、八段锦等）
    \item \imp{有害健康的风俗}：某些地区的嚼槟榔习惯（口腔癌风险）、高盐腌制食品消费（高血压、胃癌风险）、女性割礼（FGC）等陈规陋习
    \item \imp{需科学引导的风俗}：坐月子、特殊饮食禁忌等需结合现代医学知识加以科学化引导
\end{enumerate}

% ----- 6.6 宗教与道德 -----
\section{宗教与道德对健康的影响}

\subsection{宗教}

\begin{enumerate}[leftmargin=*]
    \item \imp{积极方面}：提供精神慰藉和心理支持→减轻焦虑和抑郁；宗教倡导健康生活方式（禁止酗酒、暴饮等）→减少危险行为；宗教团体提供社会网络和支持；宗教教义中的积极心理暗示
    \item \imp{消极方面}：某些宗教习俗可能延误医疗救治（拒绝输血、拒绝接种疫苗等）；宗教狂热可能导致社会冲突和暴力；特定宗教规定可能影响营养和身体健康
\end{enumerate}

\subsection{道德}

道德规范通过调节人们的社会行为间接影响健康：①促进利他行为和互助→增强社区凝聚力和社会支持；②约束危险行为（如对吸毒、酒驾的道德谴责）；③医德伦理直接关系到医疗质量和患者健康。

% ----- 6.7 非主流文化 -----
\section{非主流文化对健康的影响}

\begin{defbox}
\textbf{非主流文化（亚文化与反文化）}：与主流文化价值观念和行为规范相区别甚至相对立的文化形态和生活方式。
\end{defbox}

\begin{enumerate}[leftmargin=*]
    \item \imp{亚文化群体}可能有特殊的健康风险：某些青年亚文化中药物滥用、危险性行为等风险较高
    \item \imp{反文化}可能与主流卫生倡导相抵触（如"反疫苗运动"）
    \item 需关注文化适应（Acculturation）对移民群体健康的影响——文化冲突和文化适应压力可能导致心理问题
\end{enumerate}

% ----- 复习题 -----
\section{复习题}

\subsection*{一、选择题}
\begin{enumerate}[leftmargin=*, itemsep=3pt]
    \item 文化影响健康的特点不包括：
    \begin{enumerate}
        \item[(A)] 潜移默化性
        \item[(B)] 短暂易变性
        \item[(C)] 广泛渗透性
        \item[(D)] 双重性
    \end{enumerate}
    \item 以下哪项不是教育影响健康的机制？
    \begin{enumerate}
        \item[(A)] 提高健康素养
        \item[(B)] 改变遗传特征
        \item[(C)] 促进健康行为
        \item[(D)] 改善职业和经济地位
    \end{enumerate}
    \item 受教育程度对以下哪一人群的健康影响最为显著？
    \begin{enumerate}
        \item[(A)] 成年男性
        \item[(B)] 母亲（对儿童健康）
        \item[(C)] 老年人
        \item[(D)] 青少年
    \end{enumerate}
\end{enumerate}

\subsection*{二、名词解释}
\begin{enumerate}[leftmargin=*, itemsep=3pt]
    \item \textbf{文化（Culture）}
    \item \textbf{文化适应（Acculturation）}
    \item \textbf{亚文化（Subculture）}
\end{enumerate}

\subsection*{三、简答题}
\begin{enumerate}[leftmargin=*, itemsep=3pt]
    \item 试述文化对人群健康影响的作用模式。
    \item 简述文化影响健康的特点。
    \item 简述教育影响健康的机制。
    \item 举例说明风俗习惯对健康的影响（正反两方面）。
\end{enumerate}

\subsection*{参考答案要点}

\subsubsection*{一、选择题}
1. (B)；2. (B)；3. (B)

\subsubsection*{二、名词解释}
\begin{enumerate}[leftmargin=*]
    \item \textbf{文化}：广义指人类创造的物质和精神财富总和；狭义（社会医学关注）指社会意识形态及与之相适应的制度和组织机构，包括思想观念、价值体系、道德规范等精神文化要素。
    \item \textbf{文化适应}：不同文化群体在持续接触过程中发生的一方或双方原有文化模式变化的现象，移民群体的文化适应对其健康有重要影响。
    \item \textbf{亚文化}：社会某一群体特有的、区别于主流文化的文化模式和行为方式。
\end{enumerate}

\subsubsection*{三、简答题}
\begin{enumerate}[leftmargin=*]
    \item 四大模式：①文化决定健康观念模式（影响疾病认知和求医行为）；②文化塑造健康行为模式（饮食、饮酒等日常行为）；③文化影响卫生服务利用模式（现代vs传统医疗选择偏好）；④文化调节心理应激反应模式（应激评价和应对策略选择）。
    \item 五大特点：潜移默化性、广泛渗透性、持久稳定性、双重性（促进和损害）、群体差异性。
    \item 六方面机制：提高健康素养、促进健康行为、增强应对能力、改善职业经济地位、减少风险行为、增强社会参与。母亲教育水平对儿童健康的独立保护效应最显著。
    \item 正面：传统饮食文化（均衡搭配）、饮茶习惯、太极拳；负面：嚼槟榔（口腔癌）、高盐腌制食品（高血压/胃癌）、女性割礼。
\end{enumerate}

\newpage
